\section{Navigation Strategies}

Organisms navigating complex environments face a fundamental trade-off between exploration (gathering information about the environment) and exploitation (using existing knowledge to achieve goals efficiently). This dilemma is particularly seen in spatial navigation, where the value of exploration must be weighed against the time and energy costs of acquiring spatial knowledge that may or may not be useful (Hills et al., 2015; Mehlhorn et al., 2015).

The study of spatial navigation has focused on two strategies:

\begin{itemize}
\item \textbf{Local heuristics:} simple, egocentric rules that guide movement based on immediate sensory input without requiring a global spatial representation (O'Keefe \& Nadel, 1978; Redish, 1999). Examples include wall-following, alternation biases (tendency to alternate left-right turns), and beacon-following. These strategies require minimal memory and can be deployed immediately but may produce suboptimal results in complex environments.

\item \textbf{Cognitive maps:} mental representations of the environment that encode spatial relationships between locations, enabling flexible path planning and shortcuts (Tolman, 1948; O'Keefe \& Nadel, 1978; Moser et al., 2008). Cognitive maps are thought to be implemented in the hippocampal formation through place cells, grid cells, and associated circuits. While enabling efficient navigation, map construction requires extensive exploration and substantial neural investment.
\end{itemize}

The study conducted by Rosenberg et al.\ (2021) provides insight into this exploration behaviour in mice. Water-deprived mice learned to navigate to the hidden water port after only $\approx 10$ reward experiences --- approximately $1000$-fold faster than conventional two-alternative forced choice (2AFC) tasks where thousands of trials are required (Burgess et al., 2017). Despite achieving reliable navigation to the water port, mice spent approximately $84\%$ of their time in the maze exploring rather than executing direct paths to the goal. This proportion was remarkably consistent across animals and persisted even in unrewarded control mice (who spent $95\%$ of time exploring), indicating strong intrinsic motivation for environmental exploration.

The authors stated that the mice often followed local navigation heuristics – alternation bias, branch bias, and forward motion bias – during their exploration. Mice demonstrated accurate homing behavior from their very first deep excursion into the maze, returning directly to the entrance from endpoints without retracing their entry path. The authors suggest this could reflect immediate acquisition of some form of cognitive spatial representation. However, the binary tree structure of the maze used in this study is such that navigating from any node to the entrance did not require memorizing the route. The branch bias – the statistical preference for turning at junctions over continuing straight – is sufficient to guide the animal back to the entrance from any location in the maze. Thus, it remains possible that the mice did not learn a specific path home, but instead relied on this local turning heuristic, which reliably produces successful homing in this particular maze structure.

\subsection{Time Allocation in Navigation}
This extension focuses on identifying the trade-off between maintaining a cognitive map and relying on navigation heuristics. It uses a time-based navigation strategy comparison to calculate how many exploitation episodes are required for the cognitive map to become advantageous. We model strategy selection as a time allocation problem: animals minimize the total time required to complete a fixed number of foraging trips, subject to the constraint that some strategies require upfront temporal investment before achieving efficiency.\\
For a number of exploitation cycles, \textit{N} (number of complete foraging cycles from entrance to water to exit, referred to as \textit{bouts}), the total time for each strategy is:

\begin{itemize}
\item Heuristic strategy:
\[
T_{\text{total}} = N \times t_{\text{heuristic}}
\]

\item Cognitive map strategy:
\[
T_{\text{total}} = T_{\text{learning}} + N \times t_{\text{map}}
\]
\end{itemize}

In the above equations, $t_{\text{heuristic}}$ and $t_{\text{map}}$ represent the average time per trip for each strategy, and $T_{\text{learning}}$ represents the total time invested in exploration to construct the cognitive map. The critical breakeven point $N^*$ occurs where these total times are equal; where strategies yield equal total time:
\[
N^* = \frac{T_{\text{learning}}}{t_{\text{heuristic}} - t_{\text{map}}}
\]

For $N < N^*$, heuristics minimize total time despite being less efficient per trip, because the learning investment cannot be amortized over a sufficiently long exploitation phase. For $N > N^*$, cognitive maps become favorable. Crucially, all quantities in this framework—training time, average trip duration, and breakeven point—are directly measured from simulation rather than assumed, providing a parameter-free comparison grounded in observable behavior.

\subsection{Computational Model}
We implemented 2 agents operating in a binary tree labyrinth (k=2, depth=6) matching the exact structure used by the authors: 63 T-junctions connected by corridors leading to 64 endpoints, one containing a water port. Additionally, we increased the complexity of the labyrinth by increasing the value of the branching factor, \textit{k} and ranging the depth (number of hierarchical levels), \textit{D}, from 3-6.

\subsubsection{Cognitive Map Agent}
Implemented as tabular Q-learning (Watkins \& Dayan, 1992), this agent stores explicit state-action values representing a computational cognitive map.

\subsubsection{Heuristic Agent}
The heuristic agent was implemented using probabilistic navigation rules extracted from Rosenberg et al.'s analysis. This agent followed a decision hierarchy (an order of priorities):
\begin{itemize}
    \item Home run: If find water, reverse and go home by.
    \item Dead end: If reach a dead end, reverse 
    \item Probabilistic exploration: Sampled from the biased action distribution.
\end{itemize}
We set up biased parameters for the heuristics in the action space:
\begin{itemize}
    \item $w_\textit{forward}$ = 2.0 (Bias toward moving ahead rather than reversing)
    \item $w_\textit{alternating}$ = 1.5 (Bias toward alternating left and right turns)
    \item $w_\textit{}$
\end{itemize}
